\chapter{Estado del arte}\label{cap:estado}

Este capítulo revisa las herramientas y soluciones existentes que abordan, total o
parcialmente, los problemas descritos en el capítulo anterior. El objetivo no es
realizar un catálogo exhaustivo, sino identificar qué resuelve cada categoría de
herramientas y qué queda sin resolver — estableciendo así el hueco que CloudPorter
pretende ocupar. A continuación se analizan las principales categorías: herramientas
de infraestructura como código, soluciones de abstracción multi-cloud, herramientas
de estimación de costes, plataformas como servicio y herramientas con componente
de inteligencia artificial.

\section{Infraestructura como código}

La \textbf{infraestructura como código} (\textit{Infrastructure as Code}, \hyperlink{iac}{IaC}) es
el paradigma que permite definir y gestionar recursos cloud mediante ficheros de
configuración versionables, en lugar de a través de consolas gráficas o procesos
manuales. Es el enfoque dominante para la gestión de infraestructura y el punto de
partida natural para abordar el problema del lock-in.

\textbf{AWS CloudFormation} \cite{awsCloudFormation} y \textbf{Azure Bicep}
\cite{azureBicep} son las herramientas IaC propias de AWS y Azure respectivamente.
Permiten describir infraestructura de forma declarativa, pero están completamente
acopladas al proveedor: una plantilla CloudFormation no puede desplegarse en Azure,
ni un fichero Bicep en AWS. Representan el extremo máximo del lock-in tecnológico
— toda la infraestructura definida con estas herramientas queda irrecuperable
fuera del ecosistema de su proveedor.

\textbf{Terraform} \cite{terraformDocs} y su fork libre \textbf{OpenTofu}
\cite{opentofuDocs} son las herramientas IaC más extendidas. Soportan múltiples
proveedores cloud a través de \textit{providers} intercambiables, lo que las
posiciona frecuentemente como solución al problema del lock-in. Sin embargo, este
posicionamiento es en gran medida un mito: aunque la herramienta es la misma,
los recursos que se definen son completamente específicos de cada proveedor. Una
instancia EC2 de AWS no es equivalente a una máquina virtual de Azure en términos
de configuración — tienen parámetros, redes y modelos de seguridad distintos.
Migrar una arquitectura Terraform de AWS a Azure implica reescribir la práctica
totalidad de los ficheros de recursos \cite{cloudAgnosticMythRepo}.

\textbf{Pulumi} \cite{pulumiDocs} adopta un enfoque diferente: en lugar de un
lenguaje de configuración propio, permite definir infraestructura en lenguajes de
programación convencionales como Python, TypeScript o Go. Esto reduce la curva de
aprendizaje para desarrolladores, pero no resuelve el problema de abstracción
entre providers — los recursos siguen siendo específicos de cada cloud.

\section{Abstracción multi-cloud}

Frente a las herramientas IaC que soportan múltiples clouds sin abstraerlos,
existe una categoría de herramientas que sí intenta proporcionar una capa de
abstracción genuina entre el usuario y el proveedor.

\textbf{Crossplane} \cite{crossplaneDocs} es la propuesta más madura en este
espacio. Extiende Kubernetes mediante \textit{Custom Resource Definitions} (CRDs)
que permiten definir recursos cloud de forma agnóstica al proveedor, delegando
la traducción a controladores específicos de cada cloud. Su enfoque es sólido
técnicamente, pero introduce una dependencia significativa: requiere un clúster
Kubernetes operativo como plano de control, lo que añade una capa de complejidad
que muchos proyectos no están en condiciones de asumir.

\textbf{MultiCloudJ} \cite{multicloudj}, desarrollado por Salesforce, ofrece
una API unificada en Java para interactuar con servicios de distintos proveedores
cloud. Su enfoque es a nivel de código, no de infraestructura: abstrae llamadas
a servicios cloud pero no gestiona el despliegue ni la definición de recursos.

En el ámbito de proyectos experimentales, el repositorio
\textbf{cloud-agnostic-terraform} \cite{cloudAgnosticMythRepo} intentó definir
una capa de abstracción sobre Terraform para hacer los despliegues agnósticos al
proveedor. El proyecto lleva sin actualizaciones desde 2019, lo que ilustra tanto
la dificultad del problema como la ausencia de soluciones consolidadas en este
espacio.

\textbf{Talis} \cite{talisDocs} es un proyecto que combinaba una API de
hipervisor con Ansible para provisionar instancias en múltiples clouds. Su
alcance era limitado — centrado en cómputo y configuración inicial, sin módulos
de estimación de costes ni gestión de red — y el repositorio fue archivado en
septiembre de 2025.

\textbf{Score} \cite{scoreDocs} es una especificación YAML agnóstica de plataforma,
aceptada como proyecto CNCF Sandbox en 2024, que permite describir cargas de trabajo
—contenedores, puertos y dependencias de servicios— de forma independiente al
entorno de ejecución. A partir de un único fichero \texttt{score.yaml}, herramientas
como \textit{score-k8s} o \textit{score-compose} generan la configuración específica
para Kubernetes, Docker Compose, Amazon ECS o Google Cloud Run. Su enfoque es
deliberadamente acotado al nivel de carga de trabajo: Score no provisiona la
infraestructura subyacente —redes, instancias de cómputo, bases de datos
gestionadas— ni incorpora estimación de costes, limitándose a resolver la
inconsistencia de configuración entre entornos de desarrollo y producción.

\section{Estimación de costes y gestión del gasto cloud}

La opacidad en la estimación de costes cloud, descrita en el capítulo anterior,
ha generado un ecosistema de herramientas específicamente orientadas a la
visibilidad y el control del gasto.

\textbf{Infracost} \cite{infracostDocs} es la herramienta de referencia para
la estimación de costes en entornos IaC. Se integra con Terraform y OpenTofu,
analiza los ficheros de configuración y genera una estimación del coste mensual
de la infraestructura antes de desplegarla. Su limitación principal es que opera
sobre código Terraform ya escrito para un proveedor concreto — no permite
comparar el coste de una misma arquitectura entre distintos clouds.

En el ámbito empresarial, herramientas como \textbf{CloudHealth}
\cite{cloudHealthDocs} ofrecen visibilidad centralizada del gasto multi-cloud,
detección de recursos huérfanos y recomendaciones de optimización. Este problema
—conocido como \hyperlink{cloud-sprawl}{\textbf{cloud sprawl}}— es especialmente relevante en organizaciones
medianas y grandes: la proliferación descontrolada de recursos no monitorizados
puede representar una fracción significativa del gasto cloud total \cite{flexera2026}.
Sin embargo, estas herramientas operan sobre infraestructura ya desplegada, no
en la fase de diseño y comparación previa al despliegue.

\section{Plataformas como servicio}

Las \textbf{plataformas como servicio} (\textit{Platform as a Service}, \hyperlink{paas}{PaaS})
representan un enfoque radicalmente diferente al problema del lock-in: en lugar
de proporcionar herramientas para gestionar infraestructura en distintos clouds,
eliminan la infraestructura como concepto visible para el usuario.

\textbf{Railway} \cite{railwayDocs}, \textbf{Render} \cite{renderDocs} y
\textbf{Fly.io} \cite{flyioDocs} son ejemplos representativos de esta categoría.
El usuario despliega su aplicación sin necesitar conocer el proveedor cloud
subyacente, sin escribir IaC y sin gestionar redes o instancias. Esta simplicidad
tiene un coste: el usuario cede el control total de la infraestructura a la
plataforma, generando una nueva forma de dependencia — en este caso, con el
propio PaaS en lugar de con el cloud provider.

Este enfoque es adecuado para aplicaciones que no requieren control fino de la
infraestructura. Para cargas de trabajo que sí lo requieren —entornos empresariales,
arquitecturas complejas, requisitos de cumplimiento normativo— la abstracción
total del PaaS resulta insuficiente.

\section{Herramientas con componente de inteligencia artificial}

La incorporación de modelos de lenguaje a las herramientas de infraestructura
es una tendencia emergente que merece atención específica.

\textbf{Infracodebase} \cite{infracodebaseDocs} es una plataforma que permite
describir infraestructura en lenguaje natural y genera automáticamente código
IaC —Terraform, CloudFormation— que cumple con las políticas de seguridad y
gobierno definidas por la organización. Soporta múltiples proveedores y se integra
con pipelines CI/CD existentes. Su diferenciador es una capa de fundación que
codifica las restricciones organizacionales, asegurando que la IA genere
infraestructura coherente con los estándares internos.

\textbf{Syncable CLI} \cite{syncableCli} es una herramienta de línea de comandos
que combina análisis de seguridad, validación de IaC y despliegue cloud con
capacidades de agente de IA, permitiendo que herramientas como Claude Code o
Cursor interactúen con la infraestructura mediante lenguaje natural.

El denominador común de estas herramientas es que delegan decisiones de
infraestructura a modelos de lenguaje sin una capa de validación determinista
final. Esto introduce riesgos en un contexto donde los errores de infraestructura
tienen consecuencias operativas y económicas directas.

\section{Análisis comparativo}

La tabla~\ref{tab:comparativa} resume las capacidades de las herramientas
analizadas frente a los problemas identificados en el capítulo anterior.

\begin{table}[H]
\centering
\begin{tabular}{lccc}
\hline
\textbf{Herramienta} & \textbf{Lock-in infra} & \textbf{Lock-in conocimiento} & \textbf{Lock-in económico} \\
\hline
CloudFormation / Bicep  & No      & No      & No      \\
Terraform / OpenTofu    & Parcial & Parcial & No      \\
Pulumi                  & Parcial & Parcial & No      \\
Crossplane              & Sí      & Sí      & No      \\
Score                   & Parcial & Parcial & No      \\
Infracost               & No      & No      & Sí      \\
CloudHealth             & No      & No      & Sí      \\
Railway / Render / Fly  & Sí      & Sí      & Parcial \\
Infracodebase           & Parcial & Sí      & No      \\
\textbf{CloudPorter}         & \textbf{Sí} & \textbf{Sí} & \textbf{Sí} \\
\hline
\end{tabular}
\caption{Comparativa de herramientas frente a las dimensiones del cloud lock-in.}
\label{tab:comparativa}
\end{table}

Ninguna de las herramientas analizadas resuelve de forma conjunta las tres
dimensiones del lock-in identificadas en el capítulo anterior. Las herramientas
IaC multi-proveedor reducen parcialmente el lock-in de infraestructura pero
trasladan al usuario la responsabilidad de conocer los recursos específicos de
cada cloud, sin abordar el lock-in de conocimiento ni el económico. Las
herramientas de abstracción multi-cloud eliminan esa responsabilidad pero
requieren infraestructura adicional como plano de control y no ofrecen
visibilidad económica. Las plataformas como servicio abstraen completamente la
infraestructura a costa de generar una nueva forma de dependencia con el propio
PaaS. Las herramientas de gestión del gasto cloud permiten estimar costes antes
del despliegue, pero operan sobre infraestructura ya definida para un proveedor
concreto sin posibilidad de hacer un analisis entre otros 
provedores. Las herramientas con componente de inteligencia artificial delegan
decisiones críticas a modelos de lenguaje sin garantías deterministas.
CloudPorter propone abordar las tres dimensiones de forma integrada, con un
enfoque determinista y sin requerir infraestructura adicional.
